\section{Fazit und Ausblick}
\label{sec:fazit}

%% TODO: Dieses Kapitel ist ein Platzhalter und sollte erst final
%% geschrieben werden, wenn Anforderungen 1, 2 und 4 (QUARC-Portierung,
%% Kraftsensor-Integration, Regelungsapplikation) umgesetzt und ausgewertet
%% sind. Der folgende Text fasst nur den Stand nach Abschluss der
%% Baseline-Messung zusammen.

Im bisherigen Projektverlauf wurde die bestehende MATLAB/KVP-Anbindung
(Friesen~\cite{friesen2021}) messtechnisch als Baseline charakterisiert
(Abschnitt~\ref{sec:baseline}). Die Ergebnisse bestätigen quantitativ, was
Friesen in seiner eigenen Bewertung bereits als \enquote{teilweise erfüllt}
einstufte: Die Zykluszeit streut trotz einer harten Untergrenze von
\SI{12}{\milli\second} bis über \SI{140}{\milli\second}, mit einer klaren
Trennung zwischen wiederkehrendem Scheduling-Jitter an Bahnsegmentübergängen
und seltenen, aber deutlich größeren \ac{tcpip}-Aussetzern. Diese Baseline
bildet den Vergleichsmaßstab für die QUARC-Portierung.

% TODO: Die weiteren Arbeitsschritte laut Aufgabenstellung sind: die Portierung der Anbindung auf eine deterministisch getaktete QUARC/Simulink-Struktur, die hard- und softwareseitige Integration eines Kraftsensors über die QUARC-I/O-Karte (unter Berücksichtigung der bei Hübel bereits genutzten ethercat-Anbindung, die Wiederholung der Latenzmessung mit dem neuen Aufbau zum direkten Vergleich mit der hier vorgestellten Baseline, sowie die Umsetzung einer totzeitkompensierten, kraftbegrenzten Regelungsapplikation.